\documentclass[a4paper,11pt]{article}
\usepackage{exam-style-842-unit}
\examinehead{842 信息化概论 · 第 8 单元练习册}{练习二 参考答案与解析}
\begin{document}

\unitanswerscover{842 信息化概论 · 单元练习册}{第 8 单元参考答案}{练习二（移动电商、支付物流与电子政务）的答案与解析；每道判断题均给出改正后的正确表述。教材页码按“扫描页−7＝印刷页”换算，引文以 OCR 原文为准。}

\noindent{\bfseries\color{ctp-mauve}【关于移动商务运营模式条数的说明】}教材 8.1.2 在“2. 移动商务的运营模式”一节中先写“以下是对这 5 种移动商务运营模式的分析”，但实际列出的标题有 6 个：内容提供商、移动门户、移动运营商、WAP 网关提供商、移动直接面向客户、移动虚拟社区。本练习册按教材实际列出的 6 类命题与作答，考试作答时以“按参与角色划分的 6 类”为准，并可补一句教材表述为“5 种”的说明，避免失分。\par\vspace{0.6em}

\answerpart{一、单项选择题（10 题，每题 3 分，共 30 分）}

\q A。\begin{analysis}教材 8.1.2 给出移动电子商务的定义：基于无线网络，运用手机、掌上电脑等移动通信设备与因特网有机结合所进行的电子商务活动，即“移动终端＋网络＋交易活动＝移动电子商务”（第 131 页）。B 是传统 PC 上网的电子商务；C 电话、传真不是移动电子商务的主要形态；D 是电子商务整合的三流，不是移动电子商务的概括公式。\end{analysis}

\q A。\begin{analysis}移动电子商务的主要特点是灵活、简单、方便，能完全根据消费者的个性化需求和喜好定制，设备选择以及提供服务与信息的方式完全由用户自己控制，用户可以随时随地获取所需的服务、应用、信息和娱乐（教材 8.1.2，第 131 页）。B、C、D 均与教材表述相反。\end{analysis}

\q D。\begin{analysis}电子支付的类型按照电子支付指令发起方式分为网上支付、电话支付、移动支付、销售点终端交易、自动柜员机交易和其他电子支付（教材 8.1.3，第 133 页）。现金支付属于传统支付方式，不是电子支付按发起方式的分类项。\end{analysis}

\q D。\begin{analysis}电子现金（E-Cash）又称数字现金，具有货币价值、可交换性、可存储性和不可重复性 4 个属性。其中不可重复性要求防止数字现金的复制和重复使用，一般的数字现金系统会建立事后检测和惩罚（教材 8.1.3，第 134 页）。A 货币价值指需有现金、银行授权信用或银行证明支持；B 可交换性指可与纸币、商品服务、账户金额等互换；C 可存储性指可存储在计算机外存、IC 卡等设备中。\end{analysis}

\q A。\begin{analysis}网络银行又被称为“3A 银行”，因为它不受时间、空间限制，能够在任何时间（Anytime）、任何地点（Anywhere）、以任何方式（Anyway）为客户提供金融服务（教材 8.1.3，第 136 页）。B 3C、C 5C、D 3P 均不是教材对网络银行的别称（“5C”在教材中指移动虚拟社区的社区、协调、商务、内容、沟通）。\end{analysis}

\q A。\begin{analysis}SET 协议保障了 Internet 上信用卡支付的安全性，利用 SET 协议制定的过程规范可以实现电子商务交易过程的机密性、认证性、数据完整性等安全要求，SET 提供商家和收单银行的认证，是目前用信用卡进行网上支付的国际标准（教材 8.1.3，第 134 页）。B FTP 是文件传输协议；C SMTP 是邮件发送协议；D TCP 是传输控制协议，均不负责信用卡支付安全。\end{analysis}

\q B。\begin{analysis}第三方支付（Third-Party Payment）狭义上是指具备一定实力和信誉保障的非银行机构，借助通信、计算机和信息安全技术，采用与各大银行签约的方式，在用户与银行支付结算系统间建立连接的电子支付模式（教材 8.1.3，第 139 页）。A 网上银行由商业银行提供；C 消费者直接转账不是第三方支付；D 政府财政支付不属于第三方支付。\end{analysis}

\q A。\begin{analysis}电子政务根据政府、企业和社会公众 3 个主体间的相互关系，分为政府与政府之间的电子政务（G2G）、政府与公众之间的电子政务（G2C）、政府与企业之间的电子政务（G2B）、政府与政府雇员之间的电子政务（G2E）4 种应用模式。G2G 包括中央政府与各级地方政府之间、不同地方政府之间、政府不同部门之间的电子政务（教材 8.2.1，第 142 页）。\end{analysis}

\q A。\begin{analysis}电子政务服务体系的灵魂和价值在于服务：通过规范公共服务方法、改进公共服务手段、提高公共服务质量，实现公共服务与信息技术发展的有机结合，从而体现电子政务服务体系对于公共服务的价值（教材 8.2.2，第 146 页）。管理、监管、审批是政府履职的手段，不是该体系“灵魂和价值”的表述。\end{analysis}

\q C。\begin{analysis}电子公务在城市管理中，公务（GBCP）由 G（政府，Government）、B（企业，Business）、C（公众和社区，Citizen and Community）、P（公共设施与公共环境，Public Facilities and Public Environment）构成，P 是城市公共管理服务物质基础的核心，eGBCP 是电子公务的模式（教材 8.2.3，第 147—148 页）。A、B、D 分别是 C、B、G 的含义。\end{analysis}

\answerpart{二、多项选择题（3 题，每题 4 分，共 12 分）}

\q ABCD。\begin{analysis}教材从参与角色分析移动商务运营模式，涉及：内容提供商、移动门户、移动运营商、WAP 网关提供商、移动直接面向客户、移动虚拟社区（教材 8.1.2，第 131—132 页）。A、B、C、D 均为其中的角色。注：教材先称“5 种”但实际列出 6 类，本册按实际列出的 6 类命题。\end{analysis}

\q ABCD。\begin{analysis}网络银行安全体系包括：服务器安全（设立防火墙、使用安全服务器、24 小时监控）、身份识别（基于 RSA 公钥密码体制的加密机制、数字签名机制和登录密码，以数字证书为唯一身份标识，设有 CA 认证机构与 CECA）、安全协议（最广泛使用的是 SSL 数据加密协议）、认证介质（密码、文件数字证书、动态口令卡、动态手机口令、移动口令牌、移动数字证书）、安全意识（教材 8.1.3，第 136—138 页）。四项全对。\end{analysis}

\q ABCD。\begin{analysis}电子政务根据政府、企业和社会公众 3 个主体间的关系分为 G2G、G2C、G2B、G2E 4 种应用模式（教材 8.2.1，第 142 页）。四项全对。\end{analysis}

\answerpart{三、判断题（6 题，每题 2 分，共 12 分）}

\q √。\begin{analysis}教材 8.1.2 指出，多数学者都将移动商务定义为电子商务的一个延伸，认为移动商务是电子商务的一个子集，并将其涵盖在电子商务的定义之下研究；也有学者把移动商务视为一个崭新的研究领域（第 130—131 页）。\end{analysis}

\q ×。改正：电子现金必须具有不可重复性，要防止数字现金的复制和重复使用，一般数字现金系统会建立事后检测和惩罚机制；可以随意复制的“电子现金”会造成重复支付和金融风险。\begin{analysis}教材 8.1.3 把不可重复性列为电子现金的基本属性之一，原因在于买方可能用同一个数字现金在不同地区、不同网上商店同时购物，造成重复使用（第 134 页）。\end{analysis}

\q √。\begin{analysis}SSL 协议由 Netscape 首先研制开发，其首要目的是在两个通信方之间提供秘密而可靠的连接，大部分 Web 服务器和浏览器都支持此协议；用户通过身份认证后，传输数据全部用会话密钥加密，且每次会话的密钥随机产生（教材 8.1.3，第 137 页）。\end{analysis}

\q ×。改正：电子政务与电子公务不是同一概念。电子政务强调政府的管理角色，以政府为中心；电子公务强调政府的服务角色，是政府、企业、公众（社区）共同参与、在生产公共产品过程中的事务。\begin{analysis}教材 8.2.3 明确区分：政务更多强调政府在管理社会和提供公共产品、公共服务中的作用，公务更多强调政府、社会组织、企业和公民在互动中为社会提供的公共产品和服务；电子公务包含以政府为中心、以非营利组织为中心、以营利组织为中心、以市民为中心 4 个方面的内容（第 147—148 页）。\end{analysis}

\q ×。改正：第三方支付机构是非银行机构，作为收、付款人的支付中介提供网络支付、预付卡、银行卡收单等服务，不能吸收公众存款，也不等同于银行。\begin{analysis}教材 8.1.3 对第三方支付的定义是“具备一定实力和信誉保障的非银行机构……在用户与银行支付结算系统间建立连接的电子支付模式”；央行 2010 年《非金融机构支付服务管理办法》也将其界定为非金融机构提供的支付服务（第 139 页）。\end{analysis}

\q √。\begin{analysis}我国电子政务服务体系的组成部分包括公共服务前台和政府办公后台两大平台：公共服务前台由政府门户网站和行政服务中心组成；政府办公后台以政府内网为平台，包括办公自动化系统、政务信息资源共享平台和协同办公集成平台（教材 8.2.2，第 146—147 页）。\end{analysis}

\answerpart{四、填空题（6 空，每空 1 分，共 6 分）}

\q 网络。\begin{analysis}教材 8.1.2 的概括公式为：移动终端＋网络＋交易活动＝移动电子商务（第 131 页）。\end{analysis}

\q Anyway。\begin{analysis}网络银行又称 3A 银行：Any time（任何时间）、Any where（任何地点）、Any way（任何方式）（教材 8.1.3，第 136 页）。\end{analysis}

\q 银行卡。\begin{analysis}国际通行的网上支付工具和支付方式主要有电子现金、银行卡支付、电子支票、电子钱包和网络银行等（教材 8.1.3，第 134 页）。\end{analysis}

\q G2B。\begin{analysis}电子政务 4 种应用模式为 G2G（政府与政府）、G2C（政府与公众）、G2B（政府与企业）、G2E（政府与政府雇员）（教材 8.2.1，第 142 页）。\end{analysis}

\q 公共服务前台。\begin{analysis}电子政务服务体系的两大平台为公共服务前台（政府门户网站＋行政服务中心）与政府办公后台（办公自动化系统＋政务信息资源共享平台＋协同办公集成平台）（教材 8.2.2，第 146—147 页）。\end{analysis}

\q 公共设施与公共环境。\begin{analysis}eGBCP 中 G＝政府、B＝企业、C＝公众和社区、P＝公共设施与公共环境；P 通过间接发送、接收信息成为城市公共管理服务物质基础的核心（教材 8.2.3，第 147—148 页）。\end{analysis}

\answerpart{五、简答题（4 题，每题 6 分，共 24 分）}

\q 参考答案要点：
（1）概念（3 分）：移动电子商务是基于无线网络，运用手机、掌上电脑等移动通信设备与因特网有机结合所进行的电子商务活动，即“移动终端＋网络＋交易活动＝移动电子商务”。它被多数学者视为电子商务的延伸与子集。移动电子商务应用非常广泛，包括移动商务、移动金融、移动广告、移动管理、移动办公、移动搜索、移动娱乐、在线游戏、移动教育、无线数据服务等。
（2）主要特点（3 分）：一是灵活、简单、方便；二是能完全根据消费者的个性化需求和喜好定制；三是设备选择以及提供服务与信息的方式完全由用户自己控制；四是用户可以随时随地获取所需的服务、应用、信息和娱乐，可以使用智能手机或 PDA 查找、选择及购买商品和服务，摆脱设备网络环境的束缚，最大限度利用商务空间。
（教材 8.1.2，第 130—131 页；对应 2021 年真题简答“移动电子商务的概念和运营模式”的前半问）

\q 参考答案要点（按参与角色划分，每类 1 分，共 6 分）：
（1）内容提供商：直接或通过移动门户间接向客户提供信息和服务，以订阅等方式收费，关键成功因素是专注专业能力、把非核心任务交给战略联盟。
（2）移动门户：通过无线网络使移动用户从不同运营商和内容提供商处获取服务，是关键的用户入口，关键成功因素是建立支持不同标准、协议和终端的灵活平台，并提供个性化、区域化服务。
（3）移动运营商：提供范围广、使用方便的业务平台，为用户提供接入并在安全、计费、支付等方面提供支持，关键成功因素是控制移动网络平台、掌握主动权（典型如中国移动的移动梦网）。
（4）WAP 网关提供商：通过 WAP 网关对 Web 服务器信息进行转换，使移动终端能够访问 Web 内容，关键成功因素是把处理功能集中在网关中，减少手机操作负载。
（5）移动直接面向客户：内容提供商通过 WAP 网关直接与移动用户进行商务活动、提供信息和服务，受移动设备屏幕与键盘限制，需解决信息可视化与操作简化问题。
（6）移动虚拟社区：为有相同爱好、经历或业务相关的用户提供交流分享场所，与协调、商务、内容、沟通并称 5C，以移动虚拟社区为基础的电子商务逐渐发展为社交化电子商务。
补充：教材还提到传统上按用户类型和市场细分分为 B2B、B2C、G2C，以及按服务特征分为语音通信、SMS 服务、移动运营商整合的 WAP 服务、移动办公服务 4 种。注：教材正文称“5 种”，实际列出 6 类，作答以实际列出的 6 类为准并加以说明。（教材 8.1.2，第 131—132 页）

\q 参考答案要点：
（1）含义（1.5 分）：电子支付是指单位或个人直接或授权他人通过电子终端发出支付指令，实现货币支付与资金转移的行为。它由消费者、厂商和金融机构使用安全电子支付手段，通过网络完成货币支付或资金流转，是电子商务系统的重要组成部分。
（2）主要类型（1.5 分）：按电子支付指令发起方式分为网上支付、电话支付、移动支付、销售点终端交易、自动柜员机交易和其他电子支付。
（3）与传统的区别（1 分）：电子支付通过数字流转完成信息传输，工作环境基于开放的互联网，使用先进通信手段，具有方便、快捷、高效、经济的优势；传统支付通过现金流转、票据转让和银行汇兑等物理实体完成。
（4）常用支付工具（2 分）：电子现金（数字现金，具有货币价值、可交换性、可存储性、不可重复性）；银行卡支付（无安全措施、通过第三方代理人、简单加密、SET 信用卡支付 4 种类型）；电子支票（主要用于企业间大额付款）；电子钱包（小额支付）；网络银行（3A 银行，两种发展模式：纯虚拟银行与传统银行延伸）。此外还有第三方支付、P2P 网贷、大数据金融、众筹、信息化金融机构、互联网金融门户等互联网金融模式。
（教材 8.1.3，第 133—140 页）

\q 参考答案要点：
（1）概念（2 分）：电子政务是指政府机关利用网络计算机及通信等信息技术，实现内部和外部管理和服务的无缝集成，使政府业务流程优化、资源整合、部门重组等，并通过政府网站将大量频繁的行政管理和日常事务按照设定的程序在网上实施，从而超越时间、空间和部门间的分隔制约，全方位地为社会以及自身提供一体化、规范、高效、优质、透明、符合惯例的管理和服务，以实现电子政务流程化、公共服务电子化、社会管理电子化和政府监管信息化等。
（2）3 个主体（1 分）：政府、企业、社会公众。
（3）4 种应用模式（3 分）：G2G——政府与政府之间、政府不同机构和部门之间，内容包括电子公文处理、电子政策法规系统、电子财政管理、政府网络业绩管理；G2C——政府与公众之间，内容包括教育培训、公众福利、电子就业、电子税务、心理健康咨询、电子医疗、电子证件服务；G2B——政府与企业之间，内容包括信息咨询服务、电子采购与招标、电子税务、电子证件办理、中小企业服务；G2E——政府与政府雇员之间，内容包括电子公文处理、电子政策法规、电子培训、电子财务管理、公务员业绩评价。
（教材 8.2.1，第 142—145 页；对应 2021、2023 年真题名词解释“电子政务”）

\answerpart{六、综合题（1 题，共 16 分）}

\q 参考答案要点：
（1）电子政务服务体系的两大平台及作用（5 分）
公共服务前台（2 分）：由政府门户网站和行政服务中心组成，是物理实体与网络虚拟相结合的服务前端。政府门户网站集中实现信息发布、政务公开、公众参与和在线办公 4 项服务职能；行政服务中心通过集中式的政府职能组织方式提供“一站式”服务，现阶段主要体现在行政审批上，减少办公地点分散、办事环节繁琐带来的不便。
政府办公后台（2 分）：以政府内网为平台，包括办公自动化系统（按事务处理层、信息管理层、决策支持层 3 个层次实现办公业务智能化、数字化、无纸化和电子化）、政务信息资源共享平台（打破部门“信息孤岛”，使信息在部门间合理流动）、协同办公集成平台（打通部门横向纵向的功能与系统集成，实现跨部门协同办公）。
二者关系（1 分）：前台面向公众提供无缝隙“一站式”服务，后台通过信息共享与协同办公为前台提供支撑，共同构成以公众为中心、以服务为灵魂、以政府内部整合为生命力的电子政务服务体系。
（2）电子政务与电子公务的区别及材料归属（4 分）
区别（2 分）：电子政务强调政府的管理角色，以政府为中心，指政府与政府、企业、公众、政府雇员之间通过电子技术实现的行政事务；电子公务强调政府的服务角色，指利用信息化手段高效实现政府、企业、公众（社区）在生产公共产品过程中的事务，包含以非营利组织、营利组织、市民为中心的活动，其模式为 eGBCP。
材料归属（2 分）：政府门户网站与行政服务中心统一受理、后台办公自动化与协同办公，属于以政府为中心的电子政务；市民通过手机 App 办理社保、公积金等“一网通办”事项属于 G2C；企业通过平台参与政府采购、税务等属于 G2B；而供水供电、公共交通等公共设施与环境（P）的协同管理，以及社会组织、企业、市民共同参与提供公共服务的活动，则属于电子公务范畴。
（3）移动支付的作用（4 分，每点 2 分）
在“一网通办”中：缴费环节实现线上支付、即时到账，使办事流程闭环，减少跑腿与现金找零，提高政务服务效率与体验；支付记录可追溯，便于对账与监管。
在农产品移动电商中：移动支付使消费者随时随地完成下单付款，是移动电子商务“移动终端＋网络＋交易活动”的关键一环；它与电子钱包、第三方支付等工具结合，降低交易成本、提升交易安全性，支撑农产品从产地直达消费者。
（4）电子商务物流与供应链管理的意义（3 分）
电子商务物流是信息化、现代化、社会化的物流，采用网络化计算机技术和现代化软硬件设备，按用户订货要求进行分类、编配、整理、配货，定时定点定量交付，可减少生产企业库存、加速资金周转、提高物流效率、降低物流成本。
供应链管理通过信息网状传递、票证单据标准化与网上传输、对消费者需求的即时反应，使农户、合作社、物流企业与消费者协同运作，缓解“丰产不丰收”、损耗高、信息不对称等问题，实现农产品上行（教材 8.1.4，第 140—141 页）。
结论（1 分）：电子政务与电子公务提升政府服务与治理能力，移动支付与智能物流打通交易与配送环节，二者共同支撑“政务服务＋农产品上行”的数字化闭环。

\end{document}
